\part{Unchecked targets and remaining gates}
\label{part:open}
\section{What has not been proved}
There is no unchecked finding offered as a theorem in this part. The
following are the unresolved targets of the task. Their state is a
current gate, not a proved negative answer. A precise counterexample to
one proposed transfer is recorded separately from a proof that no other
transfer can exist.

\begin{description}
\item[O-SH-B: density after the strict-index refutation.]
The former narrow PH2/SH(a) gate is closed by P-F57.
For the ordinary cubic-inner family, $p(f_n)>2^{cn}$ and
$s(f_n)\ge2^{cn}$ for every fixed $0<c<1/4$ eventually,
while $h=\binom{n+2}{3}^2+n+1=\Theta(n^6)$.
The fully explicit version holds for every
$n\ge2\cdot10^{80}$ with $c=1/8$.
Together with P-F31, this refutes both PH3 clauses and both
literal polynomial strict-index bounds for ordinary plethysm.
The proof uses the actual primitive dyadic Hilbert poles,
half-period subtraction, a positive normal-weight product,
and a uniform two-child increment contraction. A large
period alone would not suffice, as P-F33 demonstrates.

The separate density assertion SH(b), GCT VI
Hypothesis~1.6.5(b), remains open here. It samples the
complete stated population of nonzero plethysm inputs.
P-F48 proves that the cubic rectangles and their natural
determinant twists have exponentially small density in
both precisely encoded source populations. The unscaled
cubic family has initial multiplicity zero.  P-F62 now gives
specified polynomial odd dilations with nonzero initial
multiplicity, every realized positive stretch, and exponential
strict indices.  They are genuine exceptions in the conditioned
population, but still lie inside P-F48's exponentially sparse
family.  A new amplification must respect the exact source
population and its bitlength.
P-F63 separately refutes the broad product-closed subgroup density
clause under finite constant-additive input specifications: a fixed
positive-initial torus factor pads every positive input to a strict
saturation exception.  This padding changes the group and embedding
and does not preserve the narrow ordinary-plethysm population.
P-F67 proves that an amplification within the narrow raw-part
bitlength population would suffice even with $O(\log N)$ length
overhead and polynomially bounded fibers.  It also proves that such
an amplification contradicts the exact-shell density formulation,
not only cumulative density.  No such narrow padding map is yet
constructed.

Actual support is a different, now closed question for this
family. P-F54 gives a specified balanced cubic with every
scalar-permitted coefficient degree $d\ge(10n)^{200}$ for
$n\ge33$; finite-dimension absorption gives an absolute $C$
with $f_n(s)>0$ for all $n\ge4$ and $s\ge Cn^{199}$.
Thus polynomial actual support coexists in the actual
plethysm family with exponential strict indices.
P-F54 also gives normal rational Gorenstein surface models
with every actual degree and arbitrarily large strict indices.

The alternative average-polynomial route in P-F34/P-F42
still has an unproved lower-coefficient estimate. The
highest fifth is controlled, and the sufficient norm
condition is equivalent to the same condition for the
remaining lower coefficients. This auxiliary estimate is
no longer needed to decide PH2 or SH(a): P-F57 bypasses
the average polynomial and every nonprimitive dyadic pole.

\item[O-KRONECKER: numerical positive forms and strict indices.]
The canonical parameter clause is now refuted on the actual
triple-rectangular Kronecker family by P-F60: its multiplication-
compatible canonical ring has no h.s.o.p. with all degrees below
$2^{n-1}$, while the source rank is $n^2+5n$.
The numerical positive-form clause and polynomial strict-index
bounds remain open.  The balanced chain tensor gives a smooth
fixed curve and an explicit local moving determinant.  To infer
an actual pole one must control every other fixed component
and global cancellation; a large canonical parameter alone
does not suffice. P-F64 constructs a second positive-dimensional
fixed component, so the proposed unique-component shortcut is false.
Its known balanced array is credited to Franks--Reichenbach and
Kravtsov; the orbit and component comparison is proved here.
P-F65 now determines the second tensor's full stabilizer, proves
that its quotient point is smooth with a one-dimensional fixed
tangent, and computes its exact local conormal determinant.
Its contribution has the same phase as the chain contribution at
the principal primitive root, so these two terms cannot cancel
each other there.  The other fixed components and their total
contribution remain to be controlled.

P-F58 proves that the direct height-three rectangular transfer
vanishes in dimensions above nine, and that the exterior duality
does not preserve a fixed-rank rowwise stretching ray.
Manivel's stabilization gives an exact ray for adjoint invariants,
not the cubic family.  P-F59 constructs positive forms and
polynomial strict-index bounds efficiently for its one-row
adjoint subfamily, even though its minimum period is exponential.
Neither that positive result nor inequalities of multiplicities
settles the remaining tensor-family estimate.  The cited
Briand--Orellana--Rosas examples already refute the old zero-index
assertions.

P-F71 gives a different exact restriction dictionary.  The
diagonal $\{D:(\det D)^3=1\}$ quotient of symmetric cubic
tensors is the finite bijective normalization of the image of
the tensor quotient, with its exact Veronese grading and a
conductor-shifted coefficient inequality.  For every $n\ge4$,
the symmetric restriction module has infinite graded projective
dimension, and the quotient of its Hilbert series by the tensor
Hilbert series is not even a Laurent polynomial.  Thus the
finite-resolution shortcut to pole transfer is refuted.
Universal surjectivity is also false: at
$n=\binom{3r-1}{2}$, $r\ge2$, the symmetric ring has
$\dim B_r=1$ but the tensor ring has $\dim R_r=0$.
The restriction image is nonnormal and its normalization
does not split as a module over that image.  Its conductor
has positive degree at least $\lceil\sqrt n\rceil$;
no effective upper bound is supplied.

\item[O-COUNT: uniform effective shifted counting.]
P-F2 constructs a single shifted polytope from explicit positive Hilbert
data. P-F12 controls their existence and size under a bounded
positive-form premise in expanded parameters. P-F31 shows that the
proposed numerical PH3 bound cannot supply that premise in general. P-F13 reconstructs
and verifies candidate data from a finite prefix under an explicit
denominator promise. A procedure constructing them with the original
partition-input complexity bounds remains to be given. A small object
whose existence is known separately for every instance is not automatically
a polynomial-time constructible object. The homogeneous version is already
refuted by P-K1; no revision should silently restore that false assertion.

\item[O-NSQ: nonstandard geometric realization.]
For the particular bases and products in GCT VII--VIII, one must construct
the category, identify the basis, control parity or shifts, establish
product and specialization compatibility, and supply size bounds.
P-F6 is a sufficient criterion after these data exist. The determinant
and theta-lift passages of the supplied note do not themselves identify
the relevant Hom spaces. P-F5 refutes the simplest attempted substitution
of ordinary interlacing for plethysm restriction.
P-F36 additionally proves that the actual cubic example has a
quadratic coefficient extension that survives every finite power
regrading. Its separate conjugate blocks cannot descend through
an ordinary rational projector, and its displayed raw coefficient
cannot be a finite graded Laurent character. A comparison must
retain that extension and its descent datum, or explicitly recover
the blocks from a conjugate sum. Arbitrary coefficient systems
and the source's relaxed minimal-polynomial conjectures remain
outside this obstruction.
P-F39 supplies an actual rational symmetrizable positive operator
for the two conjugate scalars, the exact source normalization of
its trace and norm, and positive Tate--Lefschetz models for those
two normalized coefficients. Before normalization its Laurent
coefficients have negative entries. P-F52 now constructs the
missing actual-algebra intertwiner for this operator: the short
central word $(Q_1-Q_2)^2+S-f^2$ gives precisely its matrix,
and an explicit positive tensor form proves positivity at every
real $q>0$. The complete relation and cell-element factorization
prove the scalar equation independently of the source's
inconsistent printed two-by-two generator matrices. Comparison
of the coefficient models with the geometric or theta objects,
their products, and the prescribed normalization remains open.
The explicit alternative normalization in (P-F52m)--(P-F52p)
shows that positive actual action and positive Lefschetz trace
and norm do not even force the source factor to divide the scale.
The source's actual spectral-interpolation rule does give an
exact positive factorization of that scale in (P-F52q)--(P-F52s).
Moreover (P-F52t)--(P-F52x) refute two literal common-object
dictionaries: the trace and norm cannot be the first two ordinary
lambda characters under any common scalar normalization, and
the prescribed normalized pair cannot be simultaneous traces of
one Frobenius and its exterior square over every finite extension.
These are precise excluded comparisons, not a general exclusion
of geometric or theta constructions.
The positive order in (P-F52y)--(P-F52ah) now supplies a finite
cohomological additive endofunctor whose Grothendieck-group action
is explicitly intertwined with the actual central operator.
Its basis is an integral bar-fixed lattice in the source's
convenient word form. The corresponding global generator fails
membership in two explicitly rescaled generator forms, and its
change to the displayed source candidate basis has a negative
coefficient. Thus full source-canonical products, the ambient
minimum-degree condition, monoidal structure and uniform tensor
degree are the remaining comparisons; the finite central
intertwiner and additive coefficient model are already constructed.
The extension (P-F52av)--(P-F52bc) gives an actual
four-dimensional source left ideal with an integral, bar-fixed
lattice on which both chosen rescaled generators $FQ_1,FQ_2$
and all their positive words have nonnegative Laurent matrices.
The four explicitly specified source vectors give the full
intertwiner, and finite additive functors realize these word
actions on Grothendieck groups. Coherent source relations,
canonical products and the other remaining structures are
still additional requirements. Both complementary pairs
$FQ_i,FP_i$ cannot be simultaneously nonnegative in one
basis where the actual generators do not commute.
The quadratic functor relations now have explicit natural
isomorphisms coherent on repeated equal generators.
The ordinary direct-sum model does not make these functors
self-adjoint, but (P-F52bg)--(P-F52bl) supplies a nonnegative
Laurent Gram matrix intertwining both adjoints and positive
definite at every positive parameter other than one. Thus
numerical adjoint compatibility is available; actual Hom
composition, adjunction and mixed-relation coherence remain open.
The ordinary graded semisimple version of that adjunction
dictionary is now refuted by (P-F52bm)--(P-F52br), even
with arbitrary rational or virtual source images and finite
power regrading. A faithful comparison with both generators
self-adjoint would force a positive Hermitian Hom form on
generic unit-circle parameters, contradicting the actual
source spectrum near a primitive twelfth root. A remaining
construction must relax at least one of those precise
requirements. The source canonical-basis and theta proposals
do not themselves impose them, so those proposals remain open.
The completed projector construction (P-F52bs)--(P-F52ce)
now supplies actual tangle complexes with the correct four-dimensional
source action, and extends to the full concrete 21-dimensional
source algebra.  It also works uniformly for the actual tensor-image
algebras of the $\mathfrak{sl}_2$ family.  Its completed Euler
comparison preserves multiplication and both regular actions;
the full Euler space is the larger Temperley--Lieb corner, not
identified with the nonstandard subalgebra.  Quadratic relations
hold by natural homotopy equivalences.  The remaining mixed
relations are proved only on Euler classes, not as coherent natural
homotopies.  The one-sided completion does not preserve the usual
bounded duality and supplies no finite pure canonical coefficient
objects.  P-F61 further refutes a regular higher-rank extension
using only whole isotypic projectors: the same adjoint summand
occurs in opposite flip parities.  A finer multiplicity-space
comparison, canonical indecomposable basis and finite purity
remain to be constructed.
There is now an additional exact limitation on this particular
completed category: (P-F52cf)--(P-F52ci) excludes every nonzero
actual integral Euler class lying in either individual conjugate
source block.  Its ordinary diagram coordinates are rational Laurent
series, whereas the source eigenvalues have leading coefficients
$(3\pm\sqrt5)/2$.  A coordinate ratio would force these eigenvalues
to be rational Laurent series.  The rational descended module and
conjugate sums remain realized; finer coefficient structures or
other categories are outside this obstruction.
P-F66 further excludes the actual split central eigenvalues, at
$q=3^e$, from being geometric Frobenius eigenvalues, even after
nonzero rational scaling, nonzero integer powers and Tate twists:
their two real conjugates have different absolute values.
This is a weight obstruction in the geometric, all-conjugates
mixed setting.  P-F69 nevertheless realizes each as a rationally
normalized $K$-linear Frobenius trace on the first cohomology of
a simple ordinary abelian surface over $\mathbb F_{q^{116}}$,
where $K=\mathbb Q(\sqrt{A(q)})$.  The factor $1/(2q^{26})$
is explicit and is not a Tate twist over that base field.
This closes the fiberwise larger-trace existence question.
Its specific base-extension compatibility is refuted: the model
at $q^2$ has positive full trace, whereas the quadratic extension
of the model at $q$ has negative trace.  Its full-trace sequence
also cannot eventually come from extensions of one fixed abelian
variety.  A different compatible parameter family and the
supplied theta operations remain unidentified.
P-F70 now realizes the actual specialized image algebra
$M_2(K)$ of the four-dimensional source module by rational
quasi-endomorphisms of the square of that abelian surface.
Its explicit cohomological intertwiner preserves both source
generators and all ordered words; the target $H^1$ contains two
copies of the actual module.  The prescribed central operator
is Rosati self-adjoint and positive.  This is a faithful map
of that image algebra, not of the full $21$-dimensional algebra,
and it does not geometrize formal $q$-grading or source
canonical-basis coefficient objects.
P-F72 additionally realizes each source number as an exact
component Frobenius trace with only the ordinary Tate twist
$(26)$ over $\mathbb F_q$.  It uses a different Honda--Tate
surface over $\mathbb F_{q^{117}}$, a $234$-dimensional Weil
restriction, and a rational correspondence projector with
$K$-rank $2^{117}$ in degree $117$.  The twisted image is
pure of weight $65$.  These exact trace models also fail
quadratic-base-extension compatibility, by their opposite full
trace signs.  Source products for this separate trace model,
formal graded canonical objects, and the supplied iterated
theta operations remain open.

\item[O-THETA: comparison with an identified invariant.]
P-F37 supplies the actual meromorphic automorphic ring of
five-variable cubics. P-F41 identifies its weight-$48$
Borcherds form with the ordinary degree-$80$ cubic discriminant,
including all products in the generated polynomial subalgebra.
The specified construction uses a regularized multiplicative
orthogonal theta lift, an $E_6$ quasi-pullback and restriction.
The next comparison must identify this named output with the
supplied iterated unitary lift or determinant, or produce an
identified larger graded collection with quantitative bounds.
The supplied note already mentions orthogonal groups; the group
name alone neither establishes nor excludes this comparison.
P-F46 identifies the exact positive input as the vector-valued
series $\Theta_{E_6}/\delta$, including the discriminant-group
gluing and principal coefficients $1,96,27,27$. Its ordinary
Petersson norm diverges. Its positive Fourier coefficient counts
grow faster than the actual $T^{(5)}$ Hilbert coefficients under
every positive affine reindexing. Thus the remaining comparison
must identify the regularized multiplicative operation or another
explicit coefficient-space construction; direct reindexing of
these input counts is refuted. Hofmann's unitary Borcherds
construction identifies the same form as a multiplicative
unitary lift of this input; the remaining comparison concerns
that exact regularized lift and the supplied iterated definite
unitary operations, not whether the output can be called unitary.
P-F51 further realizes the exact positive $E_6$ factor as an
ordinary convergent theta integral from a specified nonnegative
function on a definite $U(3)$ quotient. Its lattice cosets,
finite level, splitting on $U(3)\times SL_2$, individual-class
measure and $y^{-3/2}$ normalization are fixed. The compatible
right unitary group is split $U(1,1)$. The remaining comparison
must match the supplied definite chains to these specified
vectors and account for division by $\delta$ and the regularized
multiplicative operation; the positive factor itself is identified.
P-F53 realizes all coefficients of that exact input by finite
disjoint unions of K3 Hilbert schemes, with even cohomology
and purity in each cohomological degree. Its completed
lambda-ring identity supplies a geometric source for the
inverse-discriminant coefficient series. Its extension
(P-F53h)--(P-F53p) also identifies the known normalized cusp
product with the completed exterior Euler character of those
spaces and recovers every input coefficient by a formal inverse.
No finite actual or virtual positive-height object gives that
same exterior character. The comparison with the supplied
iterated operation and with all actual multiplicity maps is
still missing; the completed character-level operation is fixed.
P-F68 excludes a fixed finite sequence of weighted definite-theta,
ordinary finite lambda, Adams and coefficient-derivative operations
from producing that unchanged input.  Its coefficients have
polynomial growth, whereas every input component has growth at
least $2^{\lfloor\sqrt{N+1}\rfloor-1}$.
For an ordinary compact-source final rank-one theta lift, fixed
finite level and finite Gaussian-polynomial vectors place the
output in exactly this finite class.  A varying construction
with polynomial coefficient mass, exponent shifts and inverse
lattice spacing requires effective rank plus polynomial degree
$\Omega(\sqrt N/\log N)$ to match the $N$th coefficient.
The supplied chains have not been identified with such bounded
data; indefinite or regularized kernels, infinite completion,
and a direct construction of the output are still distinct
possibilities.
P-F55 identifies the square of the same orthogonal theta output
with the six-variable cubic discriminant, of coefficient degree
$192$, and matches its cyclic restriction with the square of
the five-variable discriminant. This adds an exact fixed-dimensional
compatibility; it does not supply the missing comparison with
the supplied iterated operations or an asymptotic family.
P-F45 proves that the full cyclic-cubic-cover eigenspace
classifying domain is a ball exactly for $n=4,5$. Beyond
this range an asymptotic construction must specify what
replaces that full-domain model; selected substructures,
special period images and other theta constructions remain open.

\item[O-OBSTRUCTION: an actual family of separating equations.]
In the no-occurrence range, prove a strict inequality between the actual
coordinate-ring multiplicities, or construct an explicit polynomial
vanishing on the determinant closure and not on the padded permanent.
The ambient plethysm and Kronecker multiplicities do not determine the
ranks of the quotient maps. P-F8 and PL-06 isolate this rank/comparison
requirement. No such asymptotic obstruction family is constructed here.

\item[O-DIR-HN: arbitrary solution fields in Nullstellensatz.]
C-GRH-01 proves only the cyclotomic-field transfer. Koiran's algorithm uses
a uniform prime-ideal/root-modulo-prime estimate for arbitrary primitive
solution fields. Rojas's DZH or GIPIT conditions are alternative sufficient
analytic inputs. The present supplied Dirichlet theorem does not state
these inputs. No formal nonimplication between the analytic conjectures
is proved.

\item[O-PIT: the derandomization required by GCT V.]
The exact algorithm or hardness hypothesis of each theorem must still be
provided. The supplied counting algorithm yields C-PIT-01 in a bounded
finite-field circuit model. It does not supply the characteristic-zero
black-box, border or strengthened symbolic determinant PIT assertions
used elsewhere. GRH does not replace those separate hypotheses.

\item[O-INTEGRAL: uniform bounds for specialization.]
B7.1 proves good-prime survival for a fixed integral separator; B7.2 gives
an exact rank test in any fixed degree and integer minors excluding bad
primes. An efficient GCT transfer needs a uniform separator-degree bound,
succinct rank data and a usable prime bound. Neither Hilbert-scheme
representability nor the pointwise specialization theorem supplies these
quantitative statements.

\item[O-COMPLEXITY: Boolean and algebraic separation.]
The ETH theorem refutes the precise exponential hypotheses treated in
C-SAT-01. Its supplied circuit corollary is superlinear for
$\mathsf E^{\mathsf{NP}}$, not superpolynomial for $\mathsf{NP}$.
The current constant-free integer Boolean--arithmetic disjunction cited
in the complexity section does not choose a branch and does not identify
the usual complex-field $\mathsf{VP}$ with its constant-free integer
variant. The stated upper bounds and conditional bridges do not settle
$\mathsf P$ versus $\mathsf{NP}$ or the target GCT obstruction.
\end{description}

\section{Approaches tried and the next discriminating steps}
The qualitative-sign route was tested by actual determinant projectors
and a virtual-character counterexample (P-F3--P-F4). Its next step is an
identity with the intended multiplicity projector together with a size
bound. The standard-branching route was tested on $\operatorname{Sym}^2$
(P-F5); a replacement must use the actual Schur-functor homomorphism.

The purity route was tested at three increasingly specific levels:
semisimple categorical structure (P-F6), pure finite schemes encoding
\#3SAT (P-F7), and the actual canonical plethysm algebra carrying the
obstructing Hilbert function, a positive multiplicative basis and pure
intersection cohomology (PLE-F7--PLE-F8). The next step is an
equivariant graded comparison theorem, not another proof that an unrelated
cohomology group is pure. P-F10 separately tests the missing quantitative
control of a grading.

The Hilbert-series route made positive progress: P-F1 constructs the
optimal shifted model on the obstructing ray; PLE-F5 realizes its normal
invariant ring; PLE-F6 classifies all rescalings that repair homogeneous
Ehrhart realization; P-F2 treats every explicit positive rational form;
P-F9 and P-F11 supply coefficient-positivity bounds; P-F12 connects the
last bound to the numerical positive-form clause of revised PH3.
P-F13 gives finite reconstruction under a denominator promise;
P-F16--P-F18 give finite fixed-embedding and fixed-dimension bounds;
P-F19 refutes polynomial growth of the canonical parameter bound.
P-F30--P-F31 subsequently close the numerical positive-form question
negatively, using a complete support theorem and a fixed-component
calculation. P-F56--P-F57 then supply the direct strict-index
argument and refute narrow PH2 and SH(a). Uniform construction
of a suitable revised effective counting model remains separate.
The ordinary pole is proved by these additional arguments, not
inferred from the h.s.o.p. bound alone.
P-F20 confirms that distinction with universal denominators and
fixed-target covariants. Its multivariate counterexample defeats a
transfer based only on quasi-polynomiality from an affine slice to
simultaneous rays. P-F21--P-F22 then defeat generic normal,
Cohen--Macaulay, rational and Gorenstein attempts to repair the
parameter-degree inference. P-F23--P-F24 made the geometric gate precise: the known chain
component is smooth, and the other fixed components had to be
controlled. P-F26 identifies the entire smooth fixed locus with
the chain orbit; P-F27 excludes every graph-support boundary
competitor, including rank-deficient supports and lifts of larger
two-power order. P-F29 handles the complementary cube-and-squarefree
case. P-F30 excludes all remaining mixed supports, and P-F31 applies
the noncancellation criterion. No competing nonzero component remains. The binary determinant subfamily is separately treated
in P-F25: its realized holes are bounded by thirteen, but the
source's strict saturation definition also tests off-residue
constituent values and is not settled by that table.

The archive route distinguishes the three products on representations,
separates the two Burnside operations, checks the transcribed identities
on small groups, and treats the modular and quadratic-form caveats
(PL-01--PL-08). Its obstruction gate is the actual rank of the
Foulkes--Howe multiplicity-space map. A Schur-positive character of its
source and target does not compute that rank.

The norm and approximation routes were tested against scaling, exact
zero detection, adjacent integer multiplicities, and an explicit rank-two
family approaching rank one in normalized norm (B4.1--B5.3). They do give
valid contraction-model obstructions under the stated operator bounds.
A full algebraic-orbit obstruction requires an invariant and a certified
separating gap that survive the allowed degenerations.

The arithmetic and ETH routes track the exact analytic family, all input
and reduction sizes, the direction of reductions, arithmetic precision,
and the class in which a lower bound is obtained. Their next steps are
specified in O-DIR-HN, O-PIT and O-COMPLEXITY. These are not resolved by
relabeling the supplied subexponential bound as polynomial time.

The exterior-volume route is now an exact Fermat-evaluation
dictionary (P-F47). A prime-characteristic determinant identity,
applied through normalized Hessians and Cayley's contraction,
gives the actual invariants of P-F49--P-F50. That point-specific
conductor does not imply a strict-index bound; P-F57 instead
proves exponential strict indices by the separate signed-wave
argument. P-F48 separately closes the proposed density
amplification by ordinary determinant twists, with the exact
rank and partition-size conventions of each source clause.
